\section{Related Works}
\label{sec:relatedworks}

%\bygxl{Follow the list below}
%1. scheduling, batch scheduling, continuous batching
%2. power management of GPU, Power management of LLM

\textbf{Pipeline Scheduling in LLMs.} 
Pipeline parallelism is widely adopted in frameworks such as FedPipe~\cite{fang2024automated}, PipeLLM~\cite{ma2023poster}, and CEBench~\cite{sun2024cebench} to address GPU memory limits that prevent single-device model hosting.
To maximize resources, systems like DeepSpeed~\cite{aminabadi2022deepspeed}, vLLM~\cite{kwon2023efficient}, and PyTorch~\cite{paszke2019pytorch} use batch scheduling.
Eliminating pipeline bubbles is key to performance optimization.
Prior works~\cite{arfeen2024pipefill, zheng2024response, lin2024syncintellects} focus on batch alignment, while micro-batching~\cite{athlur2022varuna, narayanan2019pipedream} enables smaller, parallelized batches.
Continuous batching, from Orca~\cite{yu2022orca}, schedules at token-level granularity to minimize misalignment bubbles and is adopted in vLLM~\cite{kwon2023efficient} and Sarathi~\cite{agrawal2024taming}.
However, we show that inherent bubbles from compute-memory mismatches cause significant energy waste.
We address this by controlling the frequency to extend computation time and reduce bubbles.

\noindent
\textbf{GPU Energy Effiency.}
Dynamic Voltage and Frequency Scaling (DVFS) dynamically adjusts CPU voltage and frequency to optimize energy efficiency, and has been extended to GPUs~\cite{mei2013measurement, tang2019impact, kakolyris2024slo, jia2024pccl}.
Research, such as~\cite{guerreiro2019dvfs}, shows DVFS achieves 16\% average energy improvements at the workload level.
% Power-capping techniques~\cite{zhao2023sustainable, patel2024characterizing, stojkovic2024towards} significantly reduce LLM inference energy, while TAPAS~\cite{stojkovic2025tapas} demonstrates power-capping prevents overheating and power failures.
% Recent work~\cite{kakolyris2025throttll} optimizes GPU energy in LLMs by dynamically adjusting frequency based on batch size and KV cache usage.
% In contrast, \thisproj{} targets energy waste from millisecond-scale pipeline bubbles, using static power-capping to extend batch execution time and eliminate bubbles.
% Its model-guided approach ensures energy efficiency without compromising throughput, orthogonal to~\cite{kakolyris2025throttll}.
% Recent work~\cite{kakolyris2025throttll} optimizes GPU energy in LLMs by dynamically adjusting frequency based on batch size and KV cache usage.
% Similarly, Perseus~\cite{chung2024reducing} dynamically scales frequencies for computations on non-critical paths during training to achieve energy savings.
% However, we observe that during inference, the duration of pipeline bubbles aligns with the overhead of dynamic GPU frequency adjustments, making it challenging to eliminate these bubbles.
% In contrast, \thisproj{} targets energy waste from these millisecond-scale bubbles through static power capping, which extends batch execution times to eliminate bubbles while preserving GPU smart power management via DVFS.
Recent work throttLL’eM~\cite{kakolyris2025throttll} examines the impact of GPU frequency and environmental factors on performance and energy consumption, employing dynamic frequency scaling to enhance energy efficiency.
In contrast, \thisproj{} addresses pipeline bubbles arising from workload imbalances in pipeline parallelism.
Perseus~\cite{chung2024reducing} dynamically scales frequencies for non-critical-path computations during training and models the energy-time relationship for each computation to inform scaling decisions.
However, we observe that pipeline bubbles during inference are considerably shorter, typically lasting several to tens of milliseconds, which aligns with the overhead of frequency adjustments.
Consequently, \thisproj{} adopts static frequency setting to avoid constant frequency adjustment and reduce energy consumption.
Furthermore, we model the relationship between iteration time and frequency during inference to guide frequency settings; unlike Perseus, our approach incorporates the effects of batch size and sequence length, yielding a more granular and precise model.






